% Introduction chapter for the Icones blueprint.
%
% Conventions used throughout the blueprint:
%   \rocq{Icones.<dir>.<File>.<name>}   link to a Rocq declaration
%   \rocqok                              the surrounding env is fully Qed-formalised
%   \mathcompok                          the env is fully formalised and uses only
%                                        mathcomp-analysis (no project-specific axioms)
%   \notready                            statement is not yet proved in Rocq
%   \uses{label1, ...}                   dep-graph dependencies
%   \label{...}                          unique label, prefixed by env type
%
% The Rocq logical name of every theory file is `Icones.<dir>.<File>`
% (see `_CoqProject`), so e.g. `cones/cone.v`'s `Cone` structure becomes
% `Icones.cones.cone.Cone`.

%%%%%%%%%%%%%%%%%%%%%%%%%%%%%%%%%%%%%%%%%%%%%%%%%%%%%%%%%%%%%%%%%%%%%%%%
\chapter*{Introduction}

This blueprint accompanies a Rocq formalisation of \emph{Integration in
Cones} by Thomas Ehrhard and Guillaume Geoffroy
(\cite{ehrhard-geoffroy-icones}). It describes the mathematics of the
formalisation in mathematical English, and links each definition,
lemma, proposition and theorem to the corresponding Rocq declaration
in the source tree under \texttt{theories/}.

\paragraph{Scope.}
The Rocq development covers paper \S2--\S9 (see \texttt{PLAN.md} at the
root of the repository):
\begin{itemize}
  \item paper \S2 --- cones, the category $\Cones$ (M1);
  \item paper \S3 --- measurable cones, paths, $\FMeas$, the
    category $\MCones$ (M2);
  \item paper \S4 --- the Pettis-style integral, integrable cones,
    cone-Fubini, completeness and well-poweredness of $\ICones$ (M3);
  \item paper \S5.1 + \S5.2 --- the concrete internal hom
    $A \lin B$ as a $\Cone$, $\MCone$, then $\ICone$ structure (M4);
  \item paper \S6 --- the substochastic-kernel category $\Skern$
    and the fully-faithful embedding $\Skern \hookrightarrow \ICones$
    (Theorem~6.5, the axiom-free regression anchor, M5);
  \item paper \S5.3--\S5.5 --- the tensor product $\otimes$ and the
    symmetric monoidal closed structure (Theorem~5.15,
    Chapter~\ref{ch:tensor}), now \emph{axiom-free};
  \item paper \S7 --- the stable cartesian closed category
    $\mathbf{SCones}$ (Theorem~7.32, Chapter~\ref{ch:stable});
  \item paper \S9 --- the exponential $\mathord{!}$ comonad
    (axiom-free) and the Seely category (Chapter~\ref{ch:exponential}),
    plus the \S9.2 fixpoints.
\end{itemize}
The framing is an \emph{axiom-free core} versus a single remaining
\emph{staged} interface: everything above is axiom-free except the
\emph{existence} of the Seely isomorphisms, mechanised modulo
\texttt{theories/axioms/seely\_interface.v} (the SAFT/Yoneda
representability content; see Chapter~\ref{ch:exponential}). Out of
scope for this iteration: paper \S8 (analytic functions, $\mathrm{ACONES}$)
and the PCS embedding (\S10).

\paragraph{How to read this blueprint.}
Each chapter mirrors a paper section (or, in some cases, a small
group of related Rocq source files). Inside each chapter, mathematical
statements are presented as \texttt{lemma} / \texttt{theorem} /
\texttt{definition} environments, each tagged with the corresponding
Rocq declaration via \verb|\rocq{...}|. The \verb|\rocqok| marker
means the surrounding environment is fully formalised and closed by
\verb|Qed| in the source tree. \verb|\uses{...}| declarations
populate the dependency graph shown in the web version of the
blueprint.

\paragraph{Status.}
All chapters are populated with the actual statements and informal
proofs, mirroring the Rocq sources. The development is \texttt{Admitted}-free
and, apart from the single staged \texttt{seely\_interface}, uses only
the three classical-logic axioms inherited from
\texttt{mathcomp-analysis}; see each chapter's closing ``Status''
section for the precise per-chapter axiom budget.
